\section{Related Work}
\label{sec:related}

\paragraph{Regenerative neural dynamics.}
Growing Neural Cellular Automata learn local update rules that develop an image from a seed and can restore a damaged pattern~\cite{nca}.
Neural Cellular Automata Manifold (NCAM) extends this construction to an autoencoding setting, using dynamic convolutions to condition cellular development on an image representation~\cite{ncam}.
ReGAE continues this investigation of learned reconstruction dynamics, with revised conditioning and a denoising objective.
Its U-Net and attention components provide interactions beyond a strictly local cellular rule.
The relevant continuity is therefore the use of learned, repeatedly applicable reconstruction operators, rather than an identical cellular architecture.

\paragraph{Denoising as generative modeling.}
Bengio et al.~\cite{dae} show how alternating corruption and learned denoising can define a generative Markov chain, and introduce walkback training to address spurious modes encountered during sampling.
Alain and Bengio~\cite{alain} connect regularized reconstruction to the score of the data density, including the small-Gaussian-corruption regime.
These results motivate interpreting denoisers as learning local structure around the data distribution.
ReGAE uses a feature-space reconstruction loss, following the broader perceptual-loss approach of Johnson et al.~\cite{perceptual}, with VGG features~\cite{vgg}, and an $L_1$ latent prediction loss.
The score identities derived for particular reconstruction objectives do not directly establish a score interpretation for these losses.

\paragraph{Diffusion and continuous-time models.}
Denoising diffusion probabilistic models connect a prescribed forward noise process to a learned reverse process~\cite{ddpm}.
DDIM constructs non-Markovian processes with the same training objective and permits deterministic sampling~\cite{ddim}.
Score-based SDEs unify several such constructions and yield a probability-flow ODE with matching time-marginal distributions under the corresponding score field~\cite{sde}.
Neural ODEs more generally parameterize a continuous-time derivative and compute trajectories using a differential-equation solver~\cite{neuralode}.
These perspectives suggest useful questions about ReGAE's trajectories, stability, and noise injection.
The implementation studied here defines a discrete denoising map without explicit noise-time conditioning or a learned continuous-time velocity.
Consequently, its recurrent application is a stochastic dynamical system, while identification with a reverse diffusion or probability-flow ODE would require additional assumptions and a derivation.

\paragraph{Autoencoding and latent generative structure.}
Diffusion Autoencoders (DiffAE) are a particularly close comparison~\cite{diffae}.
They combine a learned semantic encoder with a conditional DDIM decoder, representing an image through a 512-dimensional nonspatial semantic vector and a spatial stochastic code.
The latter retains image-specific variations needed for near-exact reconstruction; the semantic vector alone is not the complete code.
DiffAE also learns a latent DDIM prior over semantic vectors for unconditional generation, using a deep multilayer perceptron with skip connections.
Latent diffusion models similarly place generative denoising in the representation space of a pretrained autoencoder~\cite{ldm}.
ReGAE's Gencoder learns relationships across a spatial array of patch representations and predicts the encodings of slightly displaced image patches.
This changes the patch support throughout the target array and couples latent denoising to spatial prediction.
Overlapping recurrent predictions then support image traversal without global coordinate conditioning.
Our focus is this spatial transition mechanism and its integration with regenerative decoding, rather than the general presence of a second latent model.

\paragraph{Learning from a single image.}
SinGAN learns a hierarchy of patch distributions at different scales from one natural image, enabling synthesis that preserves appearance and larger structures~\cite{singan}.
It establishes the importance of multiscale internal image statistics for this setting.
ReGAE studies a related division between local appearance and spatial organization through encoded patches and displaced latent prediction.
The single-image setting makes the distinction between reconstructing a particular arrangement and generating configurations with familiar local content especially visible.

\paragraph{Representation and compression.}
Minimum description length evaluates models through the length of their description together with the data description they support~\cite{mdl}.
COIN provides an operational image-compression example by fitting a coordinate network and storing its quantized weights~\cite{coin}.
Recent work on language-model compression also connects predictive distributions to explicit coding procedures~\cite{lmcompression}.
These perspectives distinguish a low-dimensional representation from a complete compressed bitstream.
ReGAE's hierarchy concerns reusable descriptions of appearance and spatial structure; assessing it as a codec would additionally require accounting for precision, model storage, residual information, and reconstruction distortion.
